\documentclass[main.tex]{subfiles}


\begin{document}

%from pagenumber to up
% \phantomsection 
% \hypertarget{toc}{}

 

 

 
   

\section{Электромагнитные колебания}


\textbf{Электромагнитные колебания}~--- это периодические изменения заряда, тока и напряжения в электрической цепи. Изучение электромагнитных колебаний удобно начинать с рассмотрения процессов в \emph{колебательном контуре}~--- замкнутой цепи из конденсатора и катушки.  

Пусть \emph{заряженный} конденсатор замыкают на катушку (рис.\,\ref{fig:p165}, слева).


    \begin{figure}[H]
    \centering
\begin{circuitikz}[circuit ee IEC,font=\small,thin,
    line cap=round,line join=round,
>={Stealth[inset=0pt,round,length=3mm, angle'=20]},
vec/.style={>={Stealth[inset=0pt,round,length=3.5mm, angle'=20]}}]

    \ctikzset{bipoles/americaninductor/coils=3}


% 1 pic

\path (0,0) --++ (2.5,0) --++ (0,1.5) coordinate (cl) --++ (-2.5,0) --++ (0,-1.5);

%L
% \path (cl) ++ (-{48*\pgflinewidth},0) ++ (0,{.5*\pgflinewidth}) coordinate (cc);
% \fill[yellow] (cc)
% arc (0:180:{13.65*\pgflinewidth})
% arc (0:180:{13.65*\pgflinewidth})
% arc (0:180:{13.65*\pgflinewidth})
% -- cycle
% ;

%C
\path[] (cl) ++ (0,-1.5) ++ (-{79*\pgflinewidth},{30*\pgflinewidth}) --++ (0,{-60*\pgflinewidth}) ++ ({-20\pgflinewidth},0) coordinate (cc);
\fill[yellow] (cc) rectangle++ ({20*\pgflinewidth},{60*\pgflinewidth});


\draw (0,0) to[*C,name=c] ++ (2.5,0) --++ (0,1.5) coordinate (cl);
\draw[line cap=butt] (cl) to[*L,mirror,fill=yellow] ++ (-2.5,0) coordinate (cr);
\draw (cr) --++ (0,-1.5);
\node[below] at (c.south) {$q_\text{max}$};
\node[left] at (c.north west) {$+$};
\node[right] at (c.north east) {$-$};


% 2 pic
\begin{scope}[xshift=3.5cm]

\path (0,0) --++ (2.5,0) --++ (0,1.5) coordinate (cl) --++ (-2.5,0) --++ (0,-1.5);

%L
\path (cl) ++ (-{48*\pgflinewidth},0) ++ (0,{.5*\pgflinewidth}) coordinate (ll);
\fill[yellow,semitransparent] (ll)
arc (0:180:{13.65*\pgflinewidth})
arc (0:180:{13.65*\pgflinewidth})
arc (0:180:{13.65*\pgflinewidth})
-- cycle
;

%C
\path[] (cl) ++ (0,-1.5) ++ (-{79*\pgflinewidth},{30*\pgflinewidth}) --++ (0,{-60*\pgflinewidth}) ++ ({-20\pgflinewidth},0) coordinate (cc);
\fill[yellow,semitransparent] (cc) rectangle++ ({20*\pgflinewidth},{60*\pgflinewidth});


\draw (0,0) to[*C,name=c] ++ (2.5,0) --++ (0,1.5) coordinate (cl);
\draw[line cap=butt] (cl) to[*L,mirror,name=l] ++ (-2.5,0) coordinate (cr);
\draw (cr) --++ (0,-1.5);
\node[below] at (c.south) {$q$};
\node[left] at (c.north west) {$+$};
\node[right] at (c.north east) {$-$};

\node[above] at ([xshift=-.0cm]l.north) {
\tikz\draw[->,red] (0,0) ++ (-.5,0) --++ (1,0) node[midway,above,black,shift={(0,9*\pgflinewidth)},inner sep=0] {$I$};
}
;

\end{scope}


% 3 pic
\begin{scope}[xshift=7cm]

\path (0,0) --++ (2.5,0) --++ (0,1.5) coordinate (cl) --++ (-2.5,0) --++ (0,-1.5);

%L
\path (cl) ++ (-{48*\pgflinewidth},0) ++ (0,{.5*\pgflinewidth}) coordinate (ll);
\fill[yellow] (ll)
arc (0:180:{13.65*\pgflinewidth})
arc (0:180:{13.65*\pgflinewidth})
arc (0:180:{13.65*\pgflinewidth})
-- cycle
;

%C
% \path[] (cl) ++ (0,-1.5) ++ (-{79*\pgflinewidth},{30*\pgflinewidth}) --++ (0,{-60*\pgflinewidth}) ++ ({-20\pgflinewidth},0) coordinate (cc);
% \fill[yellow,semitransparent] (cc) rectangle++ ({20*\pgflinewidth},{60*\pgflinewidth});


\draw (0,0) to[*C,name=c] ++ (2.5,0) --++ (0,1.5) coordinate (cl);
\draw[line cap=butt] (cl) to[*L,mirror,name=l] ++ (-2.5,0) coordinate (cr); 
\draw (cr) --++ (0,-1.5);
\node[below] at (c.south) {};
\node[left] at (c.north west) {};
\node[right] at (c.north east) {};

\node[above] at ([xshift=-.0cm]l.north) {
\tikz\draw[->,red] (0,0) ++ (-.5,0) --++ (1,0) node[midway,above,black,,shift={(0,9*\pgflinewidth)},inner sep=0] {$I_\text{max}$};
};

\end{scope}




\end{circuitikz}
\caption{Электромагнитные колебания в колебательном контуре}
\label{fig:p165}
\end{figure}


Сразу после замыкания цепи начинаются электромагнитные колебания~--- периодические изменения заряда на конденсаторе и тока в катушке (это \emph{свободные} колебания; они совершаются только за счет энергии, запасенной в контуре~--- внешних воздействий нет). Период колебаний обозначается через~$T$. При $t=0$ (рис.\,\ref{fig:p165}, слева) конденсатор с начальным зарядо~$q_\text{max}$ начинает разряжаться, ток в катушке равен нулю (ток в ней не может измениться скачком). При $0<t<T/4$ (рис.\,\ref{fig:p165}, посередине) конденсатор с зарядом~$q$ (при этом $q<q_\text{max}$) разряжается, ток~$I$ в катушке нарастает. При $T/4$ (рис.\,\ref{fig:p165}, справа) конденсатор \emph{разряжен}, ток в катушке достиг максимального значения~$I_\text{max}$.

Теперь конденсатор начинает \emph{перезаряжаться} (рис.\,\ref{fig:p166}, слева).

    \begin{figure}[H]
    \centering
\begin{circuitikz}[circuit ee IEC,font=\small,thin,
    line cap=round,line join=round,
>={Stealth[inset=0pt,round,length=3mm, angle'=20]},
vec/.style={>={Stealth[inset=0pt,round,length=3.5mm, angle'=20]}}]

    \ctikzset{bipoles/americaninductor/coils=3}


% 2 pic
\begin{scope}[xshift=3.5cm]

\path (0,0) --++ (2.5,0) --++ (0,1.5) coordinate (cl) --++ (-2.5,0) --++ (0,-1.5);

%L
\path (cl) ++ (-{48*\pgflinewidth},0) ++ (0,{.5*\pgflinewidth}) coordinate (ll);
\fill[yellow,semitransparent] (ll)
arc (0:180:{13.65*\pgflinewidth})
arc (0:180:{13.65*\pgflinewidth})
arc (0:180:{13.65*\pgflinewidth})
-- cycle
;

%C
\path[] (cl) ++ (0,-1.5) ++ (-{79*\pgflinewidth},{30*\pgflinewidth}) --++ (0,{-60*\pgflinewidth}) ++ ({-20\pgflinewidth},0) coordinate (cc);
\fill[yellow,semitransparent] (cc) rectangle++ ({20*\pgflinewidth},{60*\pgflinewidth});


\draw (0,0) to[*C,name=c] ++ (2.5,0) --++ (0,1.5) coordinate (cl);
\draw[line cap=butt] (cl) to[*L,mirror,name=l] ++ (-2.5,0) coordinate (cr); 
\draw (cr) --++ (0,-1.5);
\node[below] at (c.south) {$q$};
\node[left] at (c.north west) {$-$};
\node[right] at (c.north east) {$+$};

\node[above] at ([xshift=-.0cm]l.north) {
\tikz\draw[->,red] (0,0) ++ (-.5,0) --++ (1,0) node[midway,above,black,shift={(0,9*\pgflinewidth)},inner sep=0] {$I$};
}
;

\end{scope}


% 3 pic
\begin{scope}[xshift=7cm]

\path (0,0) --++ (2.5,0) --++ (0,1.5) coordinate (cl) --++ (-2.5,0) --++ (0,-1.5);

%L
% \path (cl) ++ (-{48*\pgflinewidth},0) ++ (0,{.5*\pgflinewidth}) coordinate (cc);
% \fill[yellow] (cc)
% arc (0:180:{13.65*\pgflinewidth})
% arc (0:180:{13.65*\pgflinewidth})
% arc (0:180:{13.65*\pgflinewidth})
% -- cycle
% ;

%C
\path[] (cl) ++ (0,-1.5) ++ (-{79*\pgflinewidth},{30*\pgflinewidth}) --++ (0,{-60*\pgflinewidth}) ++ ({-20\pgflinewidth},0) coordinate (cc);
\fill[yellow] (cc) rectangle++ ({20*\pgflinewidth},{60*\pgflinewidth});


\draw (0,0) to[*C,name=c] ++ (2.5,0) --++ (0,1.5) coordinate (cl); 
\draw[line cap=butt] (cl) to[*L,mirror,fill=yellow] ++ (-2.5,0) coordinate (cr); 
\draw (cr) --++ (0,-1.5);
\node[below] at (c.south) {$q_\text{max}$};
\node[left] at (c.north west) {$-$};
\node[right] at (c.north east) {$+$};

\end{scope}




\end{circuitikz}
\caption{Электромагнитные колебания в колебательном контуре}
\label{fig:p166}
\end{figure}


При $T/4<t<T/2$ (рис.\,\ref{fig:p166}, слева) на обкладках конденсатора появляются заряды противоположного знака по сравнению с теми, что были вначале; ток в катушке убывает. При $t=T/2$ (рис.\,\ref{fig:p166}, справа) конденсатор перезарядился~--- его заряд снова равен~$q_\text{max}$ (однако полярность другая); тока в катушке нет. 

При $T/2<t<T$ процессы в колебательном контуре повторяются, но идут в обратном направлении (<<обратно>> заряженный конденсатор разряжается; в катушке нарастает ток в противоположном направлении по сравнению с тем, что было при $0<t<T/2$).

Желтым цветом обозначено сосредоточение \emph{энергии} в колебательном контуре при колебаниях~--- элементы как бы обмениваются энергией между собой.



















 
\end{document}